\exercisesection{1}

\begin{enumerate}
\def\labelenumi{\arabic{enumi}.}
\tightlist
\item
  设 K 为域；对于 K 的每个子环 A，令 L(A) 表示 A 的局部环 \(A_{p}\) 所成的集合，其中 p 遍历 A 的素理想（把这些局部环视为 K 的子环）。
\end{enumerate}

\begin{enumerate}
\def\labelenumi{(\alph{enumi})}
\item
  K 的局部子环 M 要支配 \(A_{p} \in L(A)\) 中的一个环，当且仅当 \(A \subset M\)；被 M 支配的局部环 \(A_{p}\) 是唯一的，并对应于素理想 \(p = m(M) \cap A\)。
\item
  设 M、N 是 K 的两个局部子环，P 是由 \(M \cup N\) 生成的 K 的子环。证明下列条件等价：
\end{enumerate}

\((\alpha)\) 存在素理想 \(\mathfrak{p}\)，属于 \(\mathbf{P}\)，使得 \(\mathfrak{m}(\mathbf{M}) = \mathfrak{p} \cap \mathbf{M}\)，\(\mathfrak{m}(\mathbf{N}) = \mathfrak{p} \cap \mathbf{N}\)。

(β) 理想 \(\mathfrak{a}\) 是 \(\mathbf{P}\) 中由 \(\mathfrak{m}(\mathbf{M}) \cup \mathfrak{m}(\mathbf{N})\) 生成的，且不同于 \(\mathbf{P}\)。

(γ) 存在 K 的一个同时支配 M 和 N 的局部子环 Q。

若这些条件满足，则称 M 和 N 相关联。

\begin{enumerate}
\def\labelenumi{(\alph{enumi})}
\setcounter{enumi}{2}
\tightlist
\item
  设 A、B 是 K 的两个子环，C 是由 A ∪ B 生成的 K 的子环。证明下列条件等价：
\end{enumerate}

\((\alpha)\) 对每个包含 A、B 的局部环 \(\mathbf{Q} \subset \mathbf{K}\)，都有 \(\mathbf{A}_{\mathfrak{p}} = \mathbf{B}_{\mathfrak{q}}\)，其中 \(\mathfrak{p} = \mathfrak{m}(\mathbf{Q}) \cap \mathbf{A}\)，\(\mathfrak{q} = \mathfrak{m}(\mathbf{Q}) \cap \mathbf{B}\)。

(β) 对 C 的每个素理想 r，都有 \(A_{p} = B_{q}\)，其中 \(p = r \cap A\)，\(q = r \cap B\)。

(γ) 若 M ∈ L(A) 且 N ∈ L(B) 相关联，则二者相同。

(δ) L(A) ∩ L(B) = L(C)。

（使用 (a) 和 (b)。）

\begin{enumerate}
\def\labelenumi{\arabic{enumi}.}
\setcounter{enumi}{1}
\item
  对于整环 A，A 是赋值环的充要条件是 A 的每个理想都是不可约的（《代数》第二章 § 2 习题 16）。
\item
  证明每个赋值环都是凝聚的（第一章 § 2 习题 12）。
\item
  设 K 为域，F 是 K 的赋值环的集合，按包含关系全序。证明 F 中各环的交是 K 的赋值环。
\item
  设 K 为域，A 是 K 的整闭子域，K 是其分式域，且 \((\mathbf{A}_{\alpha})\) 是 K 的一个赋值环族，其交为 A。若 L 是 K 的扩域，则 A 在 L 中的整闭包等于 L 中支配某个 \(A_{\alpha}\) 的赋值环之交。
\end{enumerate}

¶ 6. 设 R 为环，A 为 R 的子环，使 S = R - A 非空且 S 中两个元素之积属于 S。

\begin{enumerate}
\def\labelenumi{(\alph{enumi})}
\item
  设 \(a, a'\) 是 \(\mathbf{A}\) 中的两个元素，\(s, s'\) 是 \(S\) 中的两个元素，且 \(sa \in \mathbf{A}\)、\(s'a' \in \mathbf{A}\)。证明 \(sa', s'a\) 中至少有一个属于 \(\mathbf{A}\)。由此证明，集合 \(p_1\) 由 \(a \in \mathbf{A}\) 组成，其中存在 \(s \in S\) 使 \(sa \in \mathbf{A}\)，且它是 \(\mathbf{A}\) 的素理想。
\item
  证明集合 \(\mathfrak{p}_0\) 由 \(a \in \mathbf{A}\) 组成，其中 \(sa \in \mathbf{A}\) 对所有 \(s \in \mathbf{S}\) 成立，且它是 \(\mathbf{R}\) 和 \(\mathbf{A}\) 的素理想；它是 \(\mathbf{R}\) 的最大且包含于 \(\mathbf{A}\) 的理想，且 \(\mathfrak{p}_0 \subset \mathfrak{p}_1\)。
\item
  集合 \(\mathfrak{n}\) 由 \(x \in \mathbb{R}\) 组成，其中存在 \(s \in \mathbb{S}\) 使 \(sx = 0\)，既是 \(\mathbb{R}\) 的理想，也是 \(\mathbb{A}\) 的理想，且 \(\mathfrak{n} \subset \mathfrak{p}_0\)。
\item
  环 \(\mathbf{A}\) 在 \(\mathbb{R}\) 中整闭。
\item
  若 \(\mathbf{R}\) 是域，\(\mathbf{A}\) 是 \(\mathbf{R}\) 的赋值环，\(p_0 = (0)\)，且 \(p_1\) 是 \(\mathbf{A}\) 的极大理想。
\end{enumerate}

¶ 7. 设 R 为环。若有序对 (A, p) 由 R 的子环 A 及 A 的素理想 p 组成，且不存在由 R 的子环 A' 及 A' 的素理想 p' 组成的有序对满足 A ⊂ A'、A' ≠ A 且 p' ∩ A = p，则称 (A, p) 在 R 中是极大的。

\begin{enumerate}
\def\labelenumi{(\alph{enumi})}
\item
  对于 R 的每个子环 B 及 B 的每个素理想 q，存在在 R 中极大的有序对 \((A, p)\)，使 \(A \supset B\) 且 \(q = p \cap B\)。
\item
  有序对 (A, p) 在 R 中极大的充要条件是：对于所有 \(s \in R - A\)，存在元素 \(c_{i} \in p\) 的有限族（\(1 \leqslant i \leqslant n\)），使元素 \(b = c_{1}s + c_{2}s^{2} + \cdots + c_{n}s^{n}\) 属于 A - p（使用《代数》第二章 § 2 第 5 节命题 11 的推论 3）。
\item
  证明若有序对 \((\mathbf{A},\mathfrak{p})\) 在 \(\mathbf{R}\) 中极大，则 \(\mathbf{R} - \mathbf{A}\) 中两个元素之积属于 \(\mathbf{R} - \mathbf{A}\)。（若 \(s,s'\) 属于 \(\mathbf{R} - \mathbf{A}\) 且 \(ss' = a\in \mathbf{A}\)，考虑满足下述条件的最小整数 \(n,m\)：存在 \(c_{i}\in \mathfrak{p}\)（\(1\leqslant i\leqslant n\)）使 \(b = \sum_{i}c_{i}s^{i}\in \mathbf{A} - \mathfrak{p}\)，以及 \(c_{j}'\in \mathfrak{p}\)（\(1\leqslant j\leqslant m\)）使
\end{enumerate}

\[
b ^ {\prime} = \sum_ {j} c _ {j} ^ {\prime} s ^ {\prime j} \in A - p;
\]

  例如设 \(n \geqslant m\)，考察 \(bb' \in A - p\)，由整数 \(n\) 的定义得到矛盾。）

\begin{enumerate}
\def\labelenumi{(\alph{enumi})}
\setcounter{enumi}{3}
\item
  在 (c) 的假设下，令 \(S = R - A\)；证明理想 \(p_1\) 是 \(A\) 中由 \(a \in A\) 构成的集合，其中存在 \(s \in S\) 使 \(sa \in A\)，并且它包含于 \(p\)（使用 (b)）。同理，理想 \(p_0\) 是 \(A\) 中由 \(a \in A\) 构成的集合，其中 \(sa \in A\) 对所有 \(s \in S\) 成立，并且它包含于 \(p\)（习题 6 (b)）。记 \(A' = A / p_0\)、\(p' = p / p_0\) 和 \(R' = R / p_0\)；有序对 \((A', p')\) 于是是整环 \(\mathbf{R}'\) 中的极大有序对，并且对于每个非零 \(a' \in \mathbf{A}'\)，存在 \(s' \in \mathbf{S}' = \mathbf{R}' - \mathbf{A}'\)，使 \(s'a' \in \mathbf{S}'\)。
\item
  设 \((\mathbf{A},\mathfrak{p})\) 为整域 \(\mathbf{R}\) 中的极大有序对，使得对于每个非零 \(a\in \mathbf{A}\)，存在 \(s\in S = R - A\) 使 \(sa\in S\)。记 \(\mathrm{S}_0 = \mathrm{S}\cup \{1\}\)，并令 \(\mathbb{R}_0\) 表示分式环 \(\mathrm{S}_0^{-1}\mathrm{R}\)；典范同态 \(\mathbb{R}\to \mathbb{R}_0\) 是单射，因而将 \(\mathbb{R}\) 识别为 \(\mathbb{R}_0\) 的子环。证明 \(\mathbb{R}_0\) 是域，并可识别为 \(\mathbb{R}\) 的分式域。令 \((\mathbf{B},\mathfrak{q})\) 为 \(\mathbb{R}_0\) 中的极大有序对，使 \(\mathbf{B}\supset \mathbf{A}\) 且 \(\mathfrak{q}\cap \mathbf{A} = \mathfrak{p}\)。证明 \(\mathbf{B}\cap \mathbf{R} = \mathbf{A}\)；换言之，\(\mathbf{A}\) 是 \(\mathbf{R}\) 与 \(\mathbb{R}_0\) 中某个赋值环的交，且 \(\mathfrak{p}\) 是 \(\mathbf{A}\) 与该赋值环的极大理想的交。
\item
\end{enumerate}

¶ 8. 给定环 A，若系数属于 A 的多项式 P(X \(_{1}\) , \ldots, X \(_{n}\) ) 具有形式 X \(^{\alpha}\) + ∑ \(_{\beta} c_{\beta}\) X \(^{\beta}\)，则称其为主导的，其中 β \textless{} α 属于乘积有序集 N \(^{n}\)，并且对所有满足 c \(_{\beta}\) ≠ 0 的 β 都如此。在环 R 中，若子环 A 对每个主导多项式 P ∈ A{[}X \(_{1}\) , \ldots, X \(_{n}\) {]}（任意 n）都满足：P(s \(_{1}\) , \ldots, s \(_{n}\) ) ≠ 0 对所有元素 s \(_{i}\) ∈ R − A（1 ≤ i ≤ n），则称 A 为旁赋值的。

\begin{enumerate}
\def\labelenumi{(\alph{enumi})}
\item
  证明若 \(\mathbf{A}\) 在 \(\mathbb{R}\) 中是旁赋值的，则 \(\mathbf{S} = \mathbf{R} - \mathbf{A}\) 中两个元素之积属于 \(\mathbf{S}\)。
\item
  设 A 是环 R 的子环，且 S = R - A 中两个元素之积属于 S；令 \(p_{0}\) 为 R 和 A 中由 \(a \in A\) 组成的素理想，其中 \(sa \in A\) 对所有 \(s \in S\) 成立（习题 6 (b)）。A 在 R 中为旁赋值的充要条件是 \(A/p_{0}\) 在 \(R/p_{0}\) 中为旁赋值的。
\item
  设 R 是环，\(h: \mathbf{R} \to \mathbf{R}'\) 是环同态。若 \(\mathbf{A}'\) 是 \(\mathbf{R}'\) 的旁赋值子环，\(\mathbf{A} = \bar{h}^{-1}(\mathbf{A}')\) 在 R 中旁赋值。由此证明：若 \((\mathbf{A}, p)\) 是 R 中的极大有序对（习题 7），则 A 在 R 中旁赋值（使用 (b) 及习题 7 (d)、(e)）。
\item
  设 \(\mathbf{R}'\) 为环，\(\mathbf{R}\) 为 \(\mathbf{R}'\) 的子环。若 \(\mathbf{A}\) 是 \(\mathbf{R}\) 的旁赋值子环，证明存在旁赋值子环 \(\mathbf{A}'\)，它属于 \(\mathbf{R}'\)，使 \(\mathbf{A}' \cap \mathbf{R} = \mathbf{A}\)。（若 \(S = R - A\) 且 \(S_0 = S \cup \{1\}\)，考察环 \(T' = S_0^{-1}R'\) 及典范同态 \(h: R' \to T'\)。令 \(\mathbf{B}\) 为 \(T'\) 中由 \(h(A) \cup (h(S))^{-1}\) 生成的子环，令 \(\mathfrak{b}\) 为 \(\mathbf{B}\) 中由 \((h(S))^{-1}\) 生成的理想。证明 \(\mathfrak{b} \neq \mathbf{B}\)；取极大有序对 \((\mathbf{A}^{\prime \prime}, \mathfrak{p}^{\prime \prime})\)，它属于 \(\mathbf{T}'\)，使 \(\mathbf{B} \subset \mathbf{A}''\) 且 \(\mathfrak{b} \subset \mathfrak{p}''\)，并令 \(\mathbf{A}' = \overline{h}^{-1}(\mathbf{A}'')\)。
\end{enumerate}

由此证明，若 R 是整环，则 R 的旁赋值子环正好是 R 与其分式域的赋值环的交。

¶ 9. (a) 设 R 为环，A 为 R 的子环，x 为 R 中元素，S 为 R 中由 \(x^{k}\)（\(k \geqslant 0\)）组成的集合。证明 x 在 A 上整当且仅当 \(x/1 \in A[1/x]\)（在环 \(S^{-1}R\) 中）。若 x 不在 A 上整，证明 \(A[1/x]\) 有一个包含 1/x 的极大理想 m，且 m 在典范同态 \(A \to A[1/x]\) 下的逆像是 A 的极大理想。

\begin{enumerate}
\def\labelenumi{(\alph{enumi})}
\setcounter{enumi}{1}
\tightlist
\item
  证明 A 在 R 中的整闭包 A' 等于包含 A 的 R 的旁赋值子环（习题 8）之交。（为证明 A' 包含于这些环中的每一个，使用习题 8 (a) 和习题 6 (d)；为证明若 \(x \in R\) 不在 A 上整，则存在 R 的一个在 R 中旁赋值的子环 B 使 \(x \notin B\)，使用 (a)、习题 7 (a) 和习题 8 (c)，仿照第 3 节定理 3 进行论证。）
\end{enumerate}

